\documentclass[11pt]{article}
% ===========================================================================
%  Shared preamble for the AP Calculus AB exam-preparation worksheets.
%  Input right after \documentclass. Files beginning with "_" are not compiled
%  on their own. Navy and gold, to match the course's lesson pages.
% ===========================================================================
\usepackage[letterpaper,top=2.3cm,bottom=2.2cm,left=1.7cm,right=1.7cm,headheight=28pt,footskip=30pt]{geometry}
\usepackage{amsmath,amssymb}
\usepackage[T1]{fontenc}
\usepackage[utf8]{inputenc}
\usepackage{lmodern}
\usepackage{xcolor}
\usepackage[most]{tcolorbox}
\usepackage{enumitem}
\usepackage{fancyhdr}
\usepackage{array}
\usepackage{tikz}
\usepackage{pgfplots}
\pgfplotsset{compat=1.18}

\definecolor{navy}{HTML}{161E54}
\definecolor{ink}{HTML}{27357E}
\definecolor{gold}{HTML}{B98A00}
\definecolor{goldbg}{HTML}{FFFDF3}
\definecolor{inkbg}{HTML}{F6F7FD}
\definecolor{rule}{HTML}{C9CFF0}
\definecolor{okgreen}{HTML}{1B7A44}
\definecolor{okbg}{HTML}{F1F8F3}

\pagestyle{fancy}
\fancyhf{}
\fancyhead[L]{\footnotesize Name: \underline{\hspace{4cm}}}
\fancyhead[C]{\textbf{\color{navy}Integration Academy}\\[-2pt]{\scriptsize Calculus AB \textperiodcentered{} AP\textsuperscript{\textregistered} Exam Preparation}}
\fancyfoot[L]{\tiny\color{black!55}\parbox[t]{0.78\textwidth}{AP\textsuperscript{\textregistered} is a registered trademark of the College Board, which was not involved in the production of, and does not endorse, this material. \textcopyright{} Integration Academy.}}
\fancyfoot[R]{\scriptsize\thepage}
\renewcommand{\headrulewidth}{1.2pt}
\renewcommand{\headrule}{\hbox to\headwidth{\color{gold}\leaders\hrule height \headrulewidth\hfill}}
\renewcommand{\footrulewidth}{0pt}
\setlength{\parindent}{0pt}
\setlength{\parskip}{4pt}

% Title block: \wstitle{code}{title}{unit line}
\newcommand{\wstitle}[3]{\begin{center}{\small\bfseries\color{gold}WORKSHEET #1}\\[3pt]{\LARGE\bfseries\color{navy}#2}\\[3pt]{\small\color{black!60}#3}\end{center}\vspace{2pt}}
% Section heading
\newcommand{\sech}[1]{\par\vspace{10pt}{\large\bfseries\color{navy}#1}\par\nobreak\vspace{2pt}{\color{rule}\hrule height0.8pt}\vspace{7pt}}
% A level tag printed beside a group of questions
\newcommand{\level}[1]{\par\vspace{6pt}{\footnotesize\bfseries\color{gold}\MakeUppercase{#1}}\par\nobreak\vspace{2pt}}

\newtcolorbox{ideabox}{colframe=ink,colback=inkbg,coltitle=white,colbacktitle=ink,fonttitle=\bfseries\small,title={Key ideas},arc=2pt,boxrule=0.8pt,left=8pt,right=8pt}
\newtcolorbox{example}[1]{colframe=rule,colback=white,coltitle=navy,colbacktitle=inkbg,fonttitle=\bfseries\small,title={#1},arc=2pt,boxrule=0.8pt,borderline west={3pt}{0pt}{ink},left=10pt,right=8pt}
\newtcolorbox{frbox}[1]{colframe=gold,colback=goldbg,coltitle=white,colbacktitle=gold,fonttitle=\bfseries\small,title={#1},arc=2pt,boxrule=0.8pt,breakable,left=8pt,right=8pt}
\newtcolorbox{solbox}{colframe=okgreen,colback=okbg,coltitle=white,colbacktitle=okgreen,fonttitle=\bfseries\small,title={Worked solutions},arc=2pt,boxrule=0.8pt,breakable,left=8pt,right=8pt}
\newtcolorbox{scorebox}[1]{colframe=navy,colback=white,coltitle=white,colbacktitle=navy,fonttitle=\bfseries\small,title={#1},arc=2pt,boxrule=0.8pt,breakable,left=8pt,right=8pt}

\newcommand{\soln}{\par\smallskip\textbf{\color{ink}Solution.}\ }
\newcommand{\result}[1]{\par\smallskip\textbf{#1}}
% blank working space of a given height
\newcommand{\work}[1]{\par\vspace{#1}}
% a point value in the margin of a scoring guide
\newcommand{\pts}[1]{\hfill{\footnotesize\bfseries\color{gold}[#1]}}

% Coordinate grid for graph questions: \begin{apgraph}{xmin}{xmax}{ymin}{ymax} ... \end{apgraph}
\newenvironment{apgraph}[4]{\begin{center}\begin{tikzpicture}\begin{axis}[
  axis lines=middle, xmin=#1, xmax=#2, ymin=#3, ymax=#4,
  grid=both, grid style={black!12}, xtick distance=1, ytick distance=1,
  ticklabel style={font=\scriptsize}, xlabel={$x$}, ylabel={$y$},
  width=10cm, height=7.2cm, samples=80, clip=false,
  every axis plot/.append style={line width=1.1pt, color=ink}]}{\end{axis}\end{tikzpicture}\end{center}}
\newcommand{\closeddot}[2]{\addplot[only marks, mark=*, mark size=2.2pt, color=ink] coordinates {(#1,#2)};}
\newcommand{\opendot}[2]{\addplot[only marks, mark=*, mark size=2.2pt, color=ink, mark options={fill=white}] coordinates {(#1,#2)};}

% Multiple-choice options: four across, or stacked when the options are long.
\newcommand{\opts}[4]{\par\vspace{3pt}\begin{tabular}{@{}p{0.235\linewidth}p{0.235\linewidth}p{0.235\linewidth}p{0.235\linewidth}@{}}(A)\ #1 & (B)\ #2 & (C)\ #3 & (D)\ #4\end{tabular}\par\vspace{9pt}}
\newcommand{\optsv}[4]{\par\vspace{3pt}\begin{tabular}{@{}l@{\ \ }p{0.9\linewidth}@{}}(A) & #1\\(B) & #2\\(C) & #3\\(D) & #4\end{tabular}\par\vspace{9pt}}

\fancyhead[R]{\footnotesize\color{black!60}Worksheet 8.3}
\begin{document}

\wstitle{8.3}{Rate In, Rate Out}{Unit 8: Applications of Integration}

\begin{ideabox}
\begin{itemize}[leftmargin=*,itemsep=2pt,topsep=2pt]
  \item The integral of a rate is an amount: $\displaystyle\int_{a}^{b}R(t)\,dt$ is the amount that entered from $t=a$ to $t=b$.
  \item Amount at time $t$ $=$ starting amount $+\displaystyle\int_{0}^{t}\big(\text{in}-\text{out}\big)\,dx$.
  \item The amount is increasing when the rate in exceeds the rate out.
  \item A greatest or least amount occurs at an endpoint or where rate in $=$ rate out. Compare the amount at every candidate.
\end{itemize}
\end{ideabox}

\sech{Worked examples}
\begin{example}{Example 1 \textperiodcentered{} The amount}
A tank holds 50 litres. Water enters at 12 litres per minute and leaves at $2t$ litres per minute, for $0\le t\le10$. How much is in the tank at $t=10$?\soln $50+\displaystyle\int_{0}^{10}(12-2t)\,dt=50+120-100$.\result{70 litres.}
\end{example}

\begin{example}{Example 2 \textperiodcentered{} The greatest amount}
When does that tank hold the most water, and how much?\soln $12=2t$ at $t=6$; the net rate changes from positive to negative there. Amount: $50+72-36$.\result{86 litres, at $t=6$.}
\end{example}

\begin{example}{Example 3 \textperiodcentered{} Increasing or decreasing}
Is the amount of water in that tank increasing or decreasing at $t=8$?\soln Net rate $12-16=-4$ litres per minute.\result{Decreasing, because the rate out exceeds the rate in.}
\end{example}

\newpage
\sech{Practice}
\emph{A calculator may be used where a decimal answer is requested. Give decimals to three places.}
\level{Core}
\begin{enumerate}[leftmargin=*,itemsep=2pt]
  \item A tank holds 40 litres. Liquid enters at 8 litres per minute and leaves at 5 litres per minute. How much is in the tank after 10 minutes?\work{2.3cm}
  \item $R(t)$ is the rate, in litres per hour, at which oil enters a tank. Explain the meaning of $\displaystyle\int_{0}^{6}R(t)\,dt$.\work{2.3cm}
  \item Sand enters a bin at $6t$ kg per minute and leaves at $t^2$ kg per minute, for $0\le t\le10$. On what interval is the amount of sand increasing?\work{2.3cm}
  \item The bin in Question 3 holds 10 kg at $t=0$. How much does it hold at $t=6$?\work{2.3cm}
\end{enumerate}

\level{Exam level}
\begin{enumerate}[leftmargin=*,itemsep=2pt,start=5]
  \item Coffee beans are loaded into a hopper at $L(t)=18e^{-0.15t}$ kg per minute and used at $U(t)=4+0.5t$ kg per minute, for $0\le t\le10$. The hopper holds 25 kg at $t=0$. How many kilograms are loaded during the 10 minutes?\work{2.8cm}
  \item Is the amount of beans in the hopper increasing or decreasing at $t=6$? Give a reason.\work{2.8cm}
  \item Find the amount of beans in the hopper at $t=10$.\work{2.8cm}
  \item Find the time at which the hopper holds the most beans, and that amount. Justify your answer.\work{2.8cm}
\end{enumerate}

\begin{frbox}{Question 9 \textperiodcentered{} Free-response practice}
Fuel enters a tank at the rate $F(t)=8+4\sin\!\left(\dfrac t2\right)$ litres per minute and is drawn off at the rate $D(t)=0.3t^2$ litres per minute, for $0\le t\le8$. The tank holds 30 litres at $t=0$.
\begin{enumerate}[label=(\alph*),leftmargin=*,itemsep=2pt]
  \item How many litres of fuel enter the tank for $0\le t\le8$?\work{2.6cm}
  \item Find the rate at which the amount of fuel in the tank is changing at $t=5$. Include units.\work{2.6cm}
  \item Find the greatest amount of fuel in the tank for $0\le t\le8$. Justify your answer.\work{2.6cm}
\end{enumerate}
\end{frbox}

\newpage
\begin{solbox}
\begin{enumerate}[leftmargin=*,itemsep=4pt]
  \item $40+10(8-5)=70$ litres.
  \item The number of litres of oil that enter the tank during the first 6 hours.
  \item $6t>t^2$ for $0<t<6$.
  \item $10+\displaystyle\int_{0}^{6}(6t-t^2)\,dt=10+108-72=46$ kg.
  \item $\displaystyle\int_{0}^{10}L(t)\,dt\approx93.224$ kg.
  \item $L(6)\approx7.318$ and $U(6)=7$. Increasing, because $L(6)>U(6)$.
  \item $25+\displaystyle\int_{0}^{10}\big(L(t)-U(t)\big)\,dt\approx53.224$ kg.
  \item $L(t)=U(t)$ at $t\approx6.201$. Amounts: 25 at $t=0$, about 63.244 at $t\approx6.201$, and 53.224 at $t=10$. The most is about 63.244 kg.
  \item \begin{enumerate}[label=(\alph*),leftmargin=*,itemsep=3pt]
    \item $\displaystyle\int_{0}^{8}F(t)\,dt\approx77.229$ litres.
    \item $F(5)-D(5)\approx2.894$ litres per minute.
    \item $F(t)=D(t)$ only at $t\approx5.587$. Amounts: 30 at $t=0$, $30+\displaystyle\int_{0}^{5.587}\big(F(t)-D(t)\big)\,dt\approx72.777$ at $t\approx5.587$, and about 56.029 at $t=8$. The greatest amount is about 72.777 litres.
  \end{enumerate}
\end{enumerate}
\end{solbox}

\end{document}
